\section{Reader Guide and Reading Routes}
\label{sec:oc13-reader-contract}

OC Core 1.3.3 is a long scientific monograph. It is not meant to be entered in
one way only. The purpose of this guide is to give each reader a stable route
through the same public object without exposing internal labels, generated
reference strings, or private process vocabulary.

\subsection{Reviewer route}

A hostile reviewer should begin with the front matter, the abstract and reader
contract, the public claim-boundary chapter, the theorem and proof integration
chapter, the claim-by-claim integrated argument, and the reviewer objections
chapter. This route answers the editorial questions that decide whether the
release is reviewable: what is claimed, what is not claimed, which proof or
evidence route supports each promoted sentence, and what would reopen the
sentence.

\subsection{Formal reader route}

A formal reader should begin with the formal model, the typed 1.3.3 foundation,
the K0 resolution foundation, the lifecycle and boundary chapters, the operator
semantics, and the theorem/proof integration chapters. The proof machinery
appendix, Lean subset, finite semantic checks, and proof sheets are then used as
audit material. This route keeps theorem labels subordinate to assumptions,
dependencies, witnesses, and counterexample boundaries.

\subsection{Domain reader route}

A domain reader should begin with the model grammar, the domain embedding
chapters, the bounded replay QA chapter, the empirical execution protocol, the
target-blind evidence tables, and the methods companion. This route explains
why physics, chemistry, biology, systems, and mathematics lanes have different
evidentiary force, and why bounded replay rows do not become broad full-scope
validation claims.

\subsection{Broad scientific reader route}

A broad scientific reader should read the introduction, model, boundary,
falsifiability, K-level, operator, lifecycle, proof, evidence, prior-art, and
limitations chapters as a continuous monograph. Appendices are available for
inspection, but the main route teaches the model before asking the reader to
audit raw evidence.

\subsection{Practical user route}

A practical user should begin with the release guide, the public claim-boundary
chapter, the practical-utility chapter, the methods companion, and the evidence
package. This route answers operational questions: what can be used, what must
be checked first, which command or artifact supports the use, and which failure
class stops interpretation.

\subsection{External interaction crosswalk}

\begin{longtable}{p{0.26\textwidth}p{0.36\textwidth}p{0.30\textwidth}}
\textbf{External task} & \textbf{Primary route} & \textbf{Evidence check if verification is needed} \\
\hline
Evaluate the release for review & Front matter; claim boundary; theorem/proof integration; reviewer objections & Proof sheets, Lean subset, finite semantic checks, and claim register \\
Audit the bounded theorem core & Formal model; typed foundation; theorem/proof integration & Proof machinery appendix, Lean certificate, finite-model report \\
Inspect empirical force & Domain evidence; bounded replay QA; methods companion & Target-blind table, validation report, negative controls, falsifiers \\
Check novelty and prior art & Prior-art and comparator chapter & Comparator register, novelty matrix, source anchors \\
Use the release practically & Practical utility; methods companion; release guide & Public evidence package, checksums, replay commands \\
Verify publication quality & Title pages, abstracts, table of contents, reader contracts, metadata & Editorial adversarial review summary and public payload suitability audit \\
\end{longtable}

\subsection{What this guide is not}

This guide does not replace the monograph, the proof sheets, the methods
companion, or the evidence package. It is a public navigation layer. If a route
in this guide conflicts with the claim register, proof sheet, finite witness,
replay table, or public metadata, the conflict is a release defect and must be
repaired before publication.
